\documentclass[aps,pra,preprint,nofootinbib]{revtex4-2}

\usepackage{amsmath,amssymb,mathtools,bm}
\usepackage{hyperref}
\hypersetup{hidelinks}

\newtheorem{lemma}{Lemma}
\newtheorem{proposition}{Proposition}
\newtheorem{theorem}{Theorem}
\newcommand{\Tr}{\operatorname{Tr}}
\newcommand{\Var}{\operatorname{Var}}
\newcommand{\supp}{\operatorname{supp}}
\newcommand{\rhozero}{\rho_0}
\newcommand{\Omegazero}{\Omega_0}
\newcommand{\Cmin}{C_{\min}}
\newcommand{\Cker}{C}
\newcommand{\Vimpl}{V_{\rm impl}}
\newcommand{\HT}{H_T}
\newcommand{\Htot}{H_{\rm tot}}

\begin{document}

\title{Supplemental Material for ``Minimum dynamical coupling cost of a prescribed rank-changing quantum-state jet''}
\author{Anonymous}
\affiliation{Anonymous affiliation}
\maketitle

\section{Exact unitary kernel-curvature identity}

First take a bounded self-adjoint tangent generator $K$. Let $\Omega(\bm\theta)=U(\bm\theta)\Omegazero U(\bm\theta)^\dagger$, with $U(0)=I$ and $\partial_jU(0)=-iK_j$. Write the one-coordinate expansion
\begin{equation}
 U(\theta)=I-i\theta K+\theta^2 W+o(\theta^2),
\end{equation}
where unitarity implies $W+W^\dagger=-K^2$. Then
\begin{align}
 \partial_\theta^2\Omega(0)
 &=2K\Omegazero K-K^2\Omegazero-\Omegazero K^2
 +2W_a\Omegazero-2\Omegazero W_a,
 \label{eq:sup-second}
\end{align}
where $W_a=(W-W^\dagger)/2$ is anti-Hermitian.

Let $P=\supp\rhozero$, $Q=I-P$, and $\Tr_E\Omegazero=\rhozero$. Since $Q\rhozero Q=0$, positivity implies
\begin{equation}
 (Q\otimes I)\Omegazero=\Omegazero(Q\otimes I)=0.
 \label{eq:sup-support}
\end{equation}
Projecting Eq.~\eqref{eq:sup-second} with $Q\otimes I$ on both sides eliminates every term except $2K\Omegazero K$. Taking the partial trace gives
\begin{equation}
 Q\partial_\theta^2\rho(0)Q
 =2\Tr_E[(Q\otimes I)K\Omegazero K(Q\otimes I)].
 \label{eq:sup-kernelidentity}
\end{equation}
This argument also shows why the identity is independent of the second derivative of the unitary curve. For an unbounded self-adjoint generator with finite baseline second moment, apply the same identity to bounded spectral truncations of $K$ and pass to the state-weighted quadratic-form limit. The energy-conserving direct-sum construction below is treated more directly by differentiating its trace-class branch series.

For several coordinates, Eq.~\eqref{eq:sup-kernelidentity} applies to each diagonal second derivative. Contracting with a positive physical parameter metric after a linear coordinate transformation gives the corresponding metric-covariant form.

\section{Variance decomposition and universal lower bound}

Set $\Pi=P\otimes I_E$ and $\bar\Pi=Q\otimes I_E$. Equation~\eqref{eq:sup-support} gives $\Omegazero=\Pi\Omegazero\Pi$. For a Hermitian $K$,
\begin{align}
 \Tr(\Omegazero K^2)
 &=\Tr(\Omegazero K\Pi K)
  +\Tr(\Omegazero K\bar\Pi K),\\
 \Tr(\Omegazero K)&=\Tr(\Omegazero\Pi K\Pi).
\end{align}
Therefore
\begin{equation}
 \Var_{\Omegazero}(K)
 =\Var_{\Omegazero}(\Pi K\Pi)
 +\Tr(\Omegazero K\bar\Pi K),
 \label{eq:sup-vardecomp}
\end{equation}
where the first term is nonnegative. Summing over coordinates and using Eq.~\eqref{eq:sup-kernelidentity},
\begin{equation}
 \Vimpl\ge\frac12\Tr\Cker.
 \label{eq:sup-lower}
\end{equation}
Equality requires the support-block action of every generator to have zero variance on the baseline. The constructions below choose it to vanish entirely.

\section{Second-order positivity and the minimum feasible curvature}

\subsection{Finite-dimensional Schur-complement argument}

Let
\begin{equation}
 \rho(\bm\theta)=\rhozero+\sum_j\theta_jD_j+\frac12\sum_{jk}\theta_j\theta_kR_{jk}+o(\|\bm\theta\|^2)
\end{equation}
be positive near the origin, with $PD_jP=QD_jQ=0$. Fix a real direction $u=(u_j)$ and define
\begin{equation}
 D_u=\sum_j u_jD_j,\qquad R_u=\sum_{jk}u_ju_kR_{jk}.
\end{equation}
For small scalar $t$, positivity of $\rho(tu)$ and the Schur complement of the support block yield
\begin{equation}
 QR_uQ\succeq2QD_uP\rhozero^+PD_uQ.
 \label{eq:sup-directionalPSD}
\end{equation}
Summing Eq.~\eqref{eq:sup-directionalPSD} over an orthonormal physical parameter basis gives
\begin{equation}
 \Cker=Q\sum_jR_{jj}Q
 \succeq2\sum_jQD_jP\rhozero^+PD_jQ=:\Cmin.
 \label{eq:sup-Cmin}
\end{equation}
The semidefinite tangent-cone structure underlying this step is standard \cite{Shapiro1997,BonnansCominettiShapiro1999}.

\subsection{Separable infinite dimension}

Let $\rhozero$ be trace class on a separable Hilbert space. Because it is compact, write its positive support spectral decomposition
\begin{equation}
 \rhozero=\sum_{n\ge1}\lambda_n|n\rangle\langle n|,
 \qquad \lambda_n>0,
 \qquad \sum_n\lambda_n=1.
\end{equation}
Let $P_m=\sum_{n=1}^m|n\rangle\langle n|$. Assume
\begin{equation}
 L_j:=QD_jP\rhozero^{-1/2}\in\mathcal S_2.
 \label{eq:sup-HS}
\end{equation}
Applying the finite-dimensional test-vector/Schur argument to $P_m\oplus Q$ gives the monotone sequence
\begin{equation}
 QR_{jj}Q\succeq2QD_jP_m(P_m\rhozero P_m)^{-1}P_mD_jQ.
\end{equation}
The right-hand side increases in the positive quadratic-form sense to $2L_jL_j^\dagger$. The Hilbert--Schmidt condition makes this limit positive trace class. Summing over $j$ gives
\begin{equation}
 \Cker\succeq\Cmin:=2\sum_jL_jL_j^\dagger.
 \label{eq:sup-CminInfinite}
\end{equation}
No bounded pseudoinverse of $\rhozero$ is required.

\section{Horizontal realization of the minimum curvature}

In finite dimension define
\begin{equation}
 X_j:=QD_jP\rhozero^+,
 \qquad
 K_j^{\rm hor}=i(X_j-X_j^\dagger).
 \label{eq:sup-horizontalK}
\end{equation}
A direct support/kernel block calculation gives
\begin{equation}
 -i[K_j^{\rm hor},\rhozero]=D_j.
\end{equation}
The support block of $K_j^{\rm hor}$ vanishes, so
\begin{align}
 \Var_{\rhozero}(K_j^{\rm hor})
 &=\Tr(QD_jP\rhozero^+PD_jQ),\\
 Q\partial_j^2\rho(0)Q
 &=2QD_jP\rhozero^+PD_jQ.
\end{align}
Consequently
\begin{equation}
 \sum_j\Var(K_j^{\rm hor})=\frac12\Tr\Cmin.
 \label{eq:sup-horizontalcost}
\end{equation}
For pure-boundary tangents the SLD-QFI satisfies
\begin{equation}
 \frac14\Tr H^{\rm SLD}=\frac12\Tr\Cmin,
\end{equation}
which is the usual first-order Bures/SLD horizontal minimum \cite{BraunsteinCaves1994}. This identity is used only as prior-art geometry connecting the first-order and second-order problems.

In infinite dimension the potentially unbounded target-only operator in Eq.~\eqref{eq:sup-horizontalK} is unnecessary. One may instead use the canonical purification
\begin{equation}
 |\Omega\rangle=\sum_n\sqrt{\lambda_n}|n\rangle_T|n\rangle_E
\end{equation}
and the finite-norm global tangent vector
\begin{equation}
 |\chi_j^{\rm hor}\rangle
 =\sum_n\frac{QD_j|n\rangle}{\sqrt{\lambda_n}}\otimes|n\rangle_E.
 \label{eq:sup-horizontalvector}
\end{equation}
Equation~\eqref{eq:sup-HS} is exactly the condition $\|\chi_j^{\rm hor}\|<\infty$. The bounded rank-two global generator
\begin{equation}
 K_j=i(|\chi_j^{\rm hor}\rangle\langle\Omega|-|\Omega\rangle\langle\chi_j^{\rm hor}|)
\end{equation}
then realizes the same reduced derivative and cost.

\section{Exact attainment for an arbitrary feasible prescribed curvature}

Let
\begin{equation}
 S=\frac12(\Cker-\Cmin)\succeq0.
\end{equation}
Choose an ancillary flag Hilbert space orthogonal to the baseline purification labels and a vector $|\eta\rangle\in Q\otimes E_{\rm flag}$ satisfying
\begin{equation}
 \Tr_{E_{\rm flag}}|\eta\rangle\langle\eta|=S.
\end{equation}
Then
\begin{equation}
 \Tr_E|\eta\rangle\langle\Omega|=0,
 \qquad
 \Tr_E|\eta\rangle\langle\chi_j^{\rm hor}|=0.
 \label{eq:sup-orthflag}
\end{equation}
Set
\begin{equation}
 |\chi_1\rangle=|\chi_1^{\rm hor}\rangle+|\eta\rangle,
 \qquad
 |\chi_j\rangle=|\chi_j^{\rm hor}\rangle\;(j>1).
\end{equation}
The first derivatives are unchanged by Eq.~\eqref{eq:sup-orthflag}, while
\begin{equation}
 2\sum_j\Tr_E[Q|\chi_j\rangle\langle\chi_j|Q]
 =\Cmin+2S=\Cker.
\end{equation}
Moreover
\begin{equation}
 \sum_j\|\chi_j\|^2
 =\frac12\Tr\Cmin+\Tr S
 =\frac12\Tr\Cker.
\end{equation}
Using rank-two generators $K_j=i(|\chi_j\rangle\langle\Omega|-|\Omega\rangle\langle\chi_j|)$ proves the upper bound matching Eq.~\eqref{eq:sup-lower}.

This construction also makes clear why the theorem concerns the metric-contracted kernel Hessian. Assigning all of $S$ to one coordinate changes the separate $R_{jk}$ tensors but leaves the prescribed contraction and the optimized quadratic cost unchanged. A complete tensor-valued second-order interpolation problem would require additional compatibility data.

\section{Strong stationarity and a countable joint energy basis}

Let $H_T$ be self-adjoint on a separable Hilbert space and let $U_t=e^{-iH_Tt/\hbar}$. Suppose the positive trace-class state $\rhozero$ is strongly stationary,
\begin{equation}
 U_t\rhozero U_t^\dagger=\rhozero\quad\text{for all }t.
\end{equation}
Each nonzero eigenspace of the compact operator $\rhozero$ is finite dimensional and invariant under every $U_t$. The restriction $U_t|_{\mathcal E_\lambda}$ is therefore a finite-dimensional strongly continuous unitary representation of $\mathbb R$. Such a representation is simultaneously diagonalizable and decomposes into characters $e^{-iEt/\hbar}$. Hence the occupied support admits a countable orthonormal basis $|n\rangle$ satisfying
\begin{equation}
 \rhozero|n\rangle=\lambda_n|n\rangle,
 \qquad
 H_T|n\rangle=E_n|n\rangle.
 \label{eq:sup-jointbasis}
\end{equation}
The statement concerns only the occupied support of the trace-class state; the ambient Hamiltonian may have arbitrary continuous spectral components elsewhere.

If $D_j$ is strongly stationary as an operator/form on the relevant domain, then $QD_j|n\rangle$ has the same target energy $E_n$. If the positive trace-class excess curvature $S$ strongly commutes with $H_T$, compactness gives an energy-adapted decomposition
\begin{equation}
 S=\sum_rs_r|q_r\rangle\langle q_r|,
 \qquad H_T|q_r\rangle=F_r|q_r\rangle.
 \label{eq:sup-Senergy}
\end{equation}
The energies $F_r$ need not be occupied by the baseline.

\section{Energy-conserving infinite-dimensional construction}

This section gives the construction that remains valid when Eq.~\eqref{eq:sup-Senergy} contains target-energy shells empty at baseline.

Choose any occupied joint eigenstate $|n_*\rangle$ with weight $\lambda_*>0$ and energy $E_*$. If $S\ne0$, select positive branch weights $w_r$ satisfying $\sum_rw_r=\lambda_*$. Replace the one baseline branch of weight $\lambda_*$ by mutually incoherent copies
\begin{equation}
 w_r|n_*,a_r\rangle\langle n_*,a_r|,
\end{equation}
with orthogonal ancilla input states. All other occupied baseline eigenstates retain one branch.

For each excess eigenmode $|q_r\rangle$ of energy $F_r$, choose nonnegative ancilla input and output energies
\begin{equation}
 a_r=\max(0,F_r-E_*),
 \qquad
 b_r=\max(0,E_*-F_r),
 \label{eq:sup-compensation}
\end{equation}
so that
\begin{equation}
 E_*+a_r=F_r+b_r.
\end{equation}
Use orthogonal ancilla states at equal energies whenever necessary. Define the normalized horizontal target vectors
\begin{equation}
 |h_{j,n}\rangle:=\frac{QD_j|n\rangle}{\lambda_n}.
\end{equation}
For the split branch set
\begin{align}
 |\chi^{\rm hor}_{j,*,r}\rangle&=|h_{j,n_*}\rangle\otimes|a_r\rangle,\\
 |\eta_r\rangle&=\sqrt{\frac{s_r}{w_r}}|q_r\rangle\otimes|b_r\rangle.
\end{align}
Input and output ancilla states are chosen orthogonal, so the excess flag has zero first-order reduced signature and zero horizontal cross term. Assign it to one coordinate,
\begin{equation}
 |\chi_{1,*,r}\rangle=|\chi^{\rm hor}_{1,*,r}\rangle+|\eta_r\rangle.
\end{equation}
Each component of $|\chi_{j,*,r}\rangle$ has the same total energy as its branch baseline state $|n_*,a_r\rangle$. For an unsplit branch the horizontal target energy is unchanged, so the same ancilla state suffices.

For every classical branch $\alpha$, let $|\Omega_\alpha\rangle$ be its pure global baseline vector and define
\begin{equation}
 K_{j,\alpha}
 =i(|\chi_{j,\alpha}\rangle\langle\Omega_\alpha|-|\Omega_\alpha\rangle\langle\chi_{j,\alpha}|).
 \label{eq:sup-branchK}
\end{equation}
Both vectors in Eq.~\eqref{eq:sup-branchK} lie in the same total-energy eigenspace. Thus $K_{j,\alpha}$ commutes with the branch total Hamiltonian. The orthogonal direct sum $K_j=\oplus_\alpha K_{j,\alpha}$ is self-adjoint on its standard direct-sum domain and strongly commutes with the semibounded total Hamiltonian
\begin{equation}
 H_{\rm tot}=H_T\otimes I+I\otimes H_E,
 \qquad H_E\ge0.
\end{equation}
The resulting blockwise unitary $U(\bm\theta)=\oplus_\alpha\exp[-i\sum_j\theta_jK_{j,\alpha}]$ conserves total energy exactly.

The reduced first derivative is exact because the split weights sum back to $\lambda_*$, while the excess flags disappear at first order. The horizontal part produces $\Cmin$. The excess contribution is
\begin{equation}
 2\sum_rw_r\Tr_E|\eta_r\rangle\langle\eta_r|
 =2\sum_rs_r|q_r\rangle\langle q_r|=2S.
\end{equation}
Therefore the prescribed curvature $\Cker$ is realized exactly. The cost is
\begin{align}
 \Vimpl
 &=\frac12\Tr\Cmin+\sum_rw_r\|\eta_r\|^2\\
 &=\frac12\Tr\Cmin+\sum_rs_r
 =\frac12\Tr\Cker.
 \label{eq:sup-energycost}
\end{align}
This repairs the tempting but false shortcut of normalizing a curvature block by the baseline weight in the same target-energy shell: spectator curvature can occupy a shell with zero baseline population.

\section{Trace-norm second-order control of the direct-sum family}

The global baseline in the preceding construction is the trace-class classical mixture
\begin{equation}
 \Omegazero=\sum_\alpha w_\alpha\omega_\alpha,
 \qquad \omega_\alpha=|\Omega_\alpha\rangle\langle\Omega_\alpha|.
\end{equation}
For one branch and bounded branch generators,
\begin{align}
 \|[K_{j,\alpha},\omega_\alpha]\|_1
 &\le2\|K_{j,\alpha}\|,\\
 \|[K_{j,\alpha},[K_{k,\alpha},\omega_\alpha]]\|_1
 &\le4\|K_{j,\alpha}\|\,\|K_{k,\alpha}\|.
 \label{eq:sup-majorants}
\end{align}
The total quadratic cost gives
\begin{equation}
 \sum_\alpha w_\alpha\sum_j\|K_{j,\alpha}\|^2<\infty.
\end{equation}
Cauchy--Schwarz therefore makes the weighted first-derivative majorants summable, and also
\begin{equation}
 \sum_\alpha w_\alpha\|K_{j,\alpha}\|\|K_{k,\alpha}\|<\infty.
\end{equation}
Equation~\eqref{eq:sup-majorants} then permits termwise differentiation of the branch series by dominated convergence in trace norm. The global and reduced families are trace-norm $C^2$ at the origin. No fourth moment $\sum w_\alpha\|K_\alpha\|^4$ is required.

The branch weights may additionally be chosen so that the baseline ancilla has finite mean energy whenever desired. Because the compensation energies in Eq.~\eqref{eq:sup-compensation} form a countable sequence, one can choose positive $w_r$ with both $\sum_rw_r=\lambda_*$ and $\sum_rw_ra_r<\infty$. The flag amplitude rescales as $w_r^{-1/2}$, leaving the state-weighted excess cost $w_r\|\eta_r\|^2=s_r$ unchanged.

\section{Autonomous temporal endpoint action and a fixed-shell benchmark}

For a clean single-gap exchange, let the prescribed kernel curvature lie in an endpoint sector on which
\begin{equation}
 G_{\rm ex}=2\hbar\nu Q.
\end{equation}
Then
\begin{equation}
 \mathcal A_{\rm ex}^{(2)}
 =\frac14\Tr(G_{\rm ex}\Cker)
 =\frac{\hbar\nu}{2}\Tr\Cker
 =\hbar\nu\,V_{\min}.
\end{equation}

A three-state fixed-total-energy benchmark makes the zero-work point explicit. Let $|L\rangle,|M\rangle,|U\rangle$ be orthonormal within one degenerate total-energy eigenspace, baseline $|M\rangle$, and choose
\begin{align}
 -iK_x|M\rangle&=\frac c2(|L\rangle+|U\rangle),\\
 -iK_y|M\rangle&=\frac{ic}{2}(|L\rangle-|U\rangle).
\end{align}
Then
\begin{equation}
 \Var(K_x)=\Var(K_y)=\frac{c^2}{2},
 \qquad V_{\rm impl}=c^2,
\end{equation}
while the total-energy distribution remains a delta function for the entire family. The endpoint action is $\hbar\nu c^2$.

\section{Boundary nonregularity and scope}

The fact that ordinary likelihood theory changes at parameter-space boundaries is classical and predates the present operator construction \cite{Chernoff1954,Shapiro1985,SelfLiang1987}. Likewise, non-faithful quantum states can exhibit distinctions between SLD-QFI and continuous Bures geometry \cite{Safranek2017}, and recent work gives explicit Bures geodesics for non-faithful states \cite{CarrascoSpehner2026}. The present theorem does not claim these boundary phenomena. Its additional datum is a prescribed positive target-kernel Hessian beyond the first-order minimum, and its conclusion is the exact dilation-coupling price of that datum under an energy-conservation constraint.

Channel purification/fibre geometry and noisy-metrology channel extensions are also established \cite{FujiwaraImai2008,EscherEtAl2011,DemkowiczKolodynskiGuta2012}. The theorem here is state-specific and local: it does not optimize a channel uniformly over all possible inputs and does not claim a new generic channel-QFI formula.

\bibliography{references}

\end{document}
